% 'nonacm' to remove any running header
\documentclass[nonacm,sigconf]{acmart}

\usepackage{cite}
\usepackage{amsmath,amsfonts}
\usepackage{algorithmic}
\usepackage{graphicx}
\usepackage{textcomp}
\usepackage{xcolor}
\usepackage{booktabs}
\usepackage{tagging}

\usepackage{sc25repro}
\pagenumbering{gobble} % no page numbers

\begin{document}

\twocolumn[%
{\begin{center}
\Huge
Appendix: Artifact Description
\end{center}}
]

\appendixAD

\section{Overview of Contributions and Artifacts}

\subsection{Paper's Main Contributions}

\begin{description}
\item[$C_1$] We propose TableMoE, a training-free multimodal lookup framework
that replaces redundant MoE expert computation with table lookup and expert
reuse in MoE MLLMs.
\item[$C_2$] We evaluate TableMoE against representative baselines under
different expert cache ratios and show that TableMoE delivers consistently
better latency across the cache-ratio range.
\item[$C_3$] We analyze the effect of the prefill recomputation ratio on
TableMoE and show the tradeoff between latency, cache-hit behavior, and
accuracy under different prefill recomputation settings.
\end{description}

\subsection{Computational Artifacts}

\begin{description}
\item[$A_1$] Artifact archive DOI: N/A at the AD stage. The current public
artifact consists of the TableMoE source repository
(\url{https://github.com/fanyalun/tablemoe-sc26.git}) and the released offline
tables published at
\url{https://huggingface.co/buckets/fanyafanya/ALUTs}.
\end{description}

\begin{center}
\begin{tabular}{cll}
\toprule
Artifact ID  &  Contributions &  Related \\
             &  Supported     &  Paper Elements \\
\midrule
$A_1$   &  $C_1$ & Tables II--III, Figures 9--12 \\
\midrule
$A_1$   &  $C_2$ & Figure 13 \\
\midrule
$A_1$   &  $C_3$ & Figures 14--15 \\
\bottomrule
\end{tabular}
\end{center}

\section{Artifact Identification}

\newartifact

\artrel

This artifact contains the implementation of TableMoE, including the offline
ALUT construction pipeline, the runtime reuse/cache/scheduling logic, and the
evaluation scripts used to measure latency and accuracy. The same artifact is
used to support all three contributions of the paper.

\artexp

This artifact is expected to reproduce the main latency and accuracy trends
reported in the paper through the repository's quick reproduction workflow.
The final artifact is
\path{perf_results/quick_reproduction/summary/quick_reproduction_report.md},
which contains three summary tables: (1) the accuracy comparison among
\texttt{transformers}, \texttt{skip}, and \texttt{tablemoe} at the default
keep-rate setting, (2) the performance comparison across different expert cache
ratios, and (3) the TableMoE comparison across different prefill
recomputation/keep-rate settings in terms of TTFT, TPOT, accuracy, and
decode-only cache hit rate. The companion manifest file
\path{perf_results/quick_reproduction/summary/quick_reproduction_manifest.json}
provides a machine-readable index of all intermediate experiment roots and the
configuration used to render the final report.

\arttime

\begin{description}
\item[Setup] 10--30 minutes. This covers cloning the repository, creating the
Python environment, and installing the dependencies.
\item[Optional offline-table preparation] Downloading the released offline
tables takes 10--20 minutes. Rebuilding the offline tables locally takes about
2--3 hours, depending on storage speed and the time required to profile and
cluster routed hidden states.
\item[Execution] The full quick reproduction wrapper is composed of four
stages: \texttt{perf}, \texttt{cache\_hit}, \texttt{acc}, and \texttt{report}.
In our environment, the \texttt{eval\_perf} stage takes about 20--30 minutes,
the \texttt{eval\_acc} stage takes about 2 hours, and the final report
generation is lightweight. The total runtime depends on whether the full
workflow is executed or a subset of stages is selected through \texttt{STEPS}.
\item[Analysis] Less than 5 minutes. The scripts write the main comparison
tables directly to the result directories.
\end{description}

\artin

\artinpart{Hardware}

All experiments are conducted on a CPU--GPU
heterogeneous platform using one NUMA node of an Intel Xeon Gold 5318Y CPU
(24 cores), 256~GB DDR4 memory, and one NVIDIA A100 PCIe GPU with 80~GB HBM,
connected via PCIe~4.0~$\times$16.

\artinpart{Software}

The software environment used for evaluation is configured as follows:
Ubuntu@22.04.4, Python@3.10.18, CUDA@12.8, NVIDIA Driver@580.95.05,
PyTorch@2.8.0+cu128, torchvision@0.23.0+cu128, Transformers@4.57.0,
FlashAttention@2.7.3, and Triton@3.4.0.

\artinpart{Datasets / Inputs}

We use two MoE MLLMs and five public multimodal benchmarks in our evaluation,
as listed below. The models and datasets can be downloaded and prepared using
the scripts provided in the repository README.
\begin{itemize}
\item Models: Qwen3-VL-30B-A3B-Instruct and DeepSeek-VL2
\item Datasets: RealWorldQA, MMBench\_DEV\_EN\_V11, AI2D\_TEST,
ScienceQA\_TEST, and POPE
\end{itemize}

\artinpart{Installation and Deployment}

We use a Conda-based Python environment as the base execution environment.
Please follow the instructions in Section 2 of the repository README to set up
the environment and install the required Python dependencies. If a prebuilt
flash-attn wheel is unavailable for the local environment, the same CUDA 12.8
toolchain must be used to build it locally.

\artcomp

The quick reproduction workflow is:
\begin{description}
\item[$T_1$] Prepare the environment and install the dependencies following
Section 2 of the repository README.
\item[$T_2$] Prepare the required inputs following Section 3.1 of the
repository README.
\item[$T_3$] Optionally prepare the offline tables by downloading the released
version or rebuilding them locally.
\item[$T_4$] Run the performance and cache-hit stages of
\texttt{run\_quick\_reproduction.sh}. These stages organize the outputs under
\path{perf_results/quick_reproduction/}.
\item[$T_5$] Run the accuracy stage of \texttt{run\_quick\_reproduction.sh}.
This stage writes the base comparison and the TableMoE keep-rate sweep under
the corresponding \path{acc_results/quick_reproduction/} subdirectories.
\item[$T_6$] Run the report stage of \texttt{run\_quick\_reproduction.sh}.
This stage renders the final markdown report and manifest.
\end{description}

The performance stage writes the cache-ratio sweep, keep-rate sweep, and
decode-only cache-hit statistics under the corresponding
\path{perf_results/quick_reproduction/} subdirectories.

The workflow dependency is as follows:
$T_1 \rightarrow T_2 \rightarrow T_3$, and $T_3 \rightarrow T_4$,
$T_3 \rightarrow T_5$, and $T_4, T_5 \rightarrow T_6$.
If the wrapper is invoked with a subset of \texttt{STEPS}, only the
corresponding subset of artifacts is generated. The complete
\path{quick_reproduction_report.md} is produced only when the \texttt{report}
step executes successfully.

\artout

The outputs of \texttt{run\_quick\_reproduction.sh} can be viewed at three
levels. First, the intermediate raw results are written under
\path{perf_results/quick_reproduction/} and
\path{acc_results/quick_reproduction/}. For example, the performance side
contains cache-ratio roots such as \path{cache_ratio_0_25/},
\path{cache_ratio_0_5/}, and \path{cache_ratio_0_75/}, keep-rate roots such as
\path{keep_rate_0_4/}, \path{keep_rate_0_6/}, and \path{keep_rate_0_8/}, and
cache-hit roots such as \path{cache_hit_keep_rate_0_4/},
\path{cache_hit_keep_rate_0_6/}, and \path{cache_hit_keep_rate_0_8/}. Each
\texttt{eval\_perf} run directory contains raw files such as
\path{run_config.json}, \path{output.json}, \path{metrics.json}, and
\path{accuracy.json}. Each cache-hit directory contains the core summary file
\path{decode_cache_hit_summary.json}, whose main fields include
\texttt{hits}, \texttt{resident\_hits}, \texttt{prefetched\_hits},
\texttt{misses}, \texttt{total}, and \texttt{hit\_rate}. On the accuracy side,
the wrapper writes roots such as \path{base_methods/},
\path{tablemoe_keep_rate_0_4/}, and \path{tablemoe_keep_rate_0_8/}, each of
which contains a model result directory and the corresponding
\path{<output_model_name>_summary.\{json,md\}} files.

Second, each evaluation root writes stage-level summary tables:
\begin{itemize}
\item \path{perf_results/quick_reproduction/*/summary/perf_table.\{json,csv,md\}}
\item \path{perf_results/quick_reproduction/*/summary/simple_accuracy_table.\{json,csv,md\}}
\item \path{acc_results/quick_reproduction/*/summary/accuracy_table.\{json,csv,md\}}
\end{itemize}

Third, the final top-level artifacts are:
\begin{itemize}
\item \path{perf_results/quick_reproduction/summary/quick_reproduction_report.md}
\item \path{perf_results/quick_reproduction/summary/quick_reproduction_manifest.json}
\end{itemize}

The markdown report is the primary human-readable artifact. The manifest is the
machine-readable index that records the experiment roots for the cache-ratio
sweeps, keep-rate sweeps, cache-hit sweeps, accuracy sweeps, and the
configuration used to render the final report.

\end{document}
